\section{Introduction}
\label{sec:intro}

There are two sellers, and each owns a single dataset.
%
There are $n \ge 1$ poor buyers and a single rich buyer.
Let $1 \le \alpha \le \beta$.
\begin{enumerate}
\item Each poor buyer has a budget of $b_1 = 1$,
    and values the two datasets at $v_{1,1} = \alpha$ and $v_{1,2} = 1$, respectively.
\item The rich buyer has a budget of $b_2 = \beta$.
    She values the two datasets at $v_{2,1} = \beta$ and $v_{2,2} = 0$.
\end{enumerate}

The sellers set prices $p$ and $q$ for their datasets.
The buyers are allowed to purchase fractional amounts of the datasets.
The rich buyer only purchases dataset 1, so she spends of her budget,
and that's how much revenue seller 1 earns from her.

Each poor buyer \emph{prefers} dataset 1 if $v_{1,1}/p > v_{1,2}/q$ (i.e., $p < \alpha q$),
prefers dataset 2 if $v_{1,1}/p < v_{1,2}/q$ (i.e., $p > \alpha q$),
and prefers both datasets equally if $v_{1,1}/p < v_{1,2}/q$ (i.e., $p = \alpha q$).
If she prefers both datasets equally, she would either fully purchase both,
or spend her entire budget on purchasing both datasets, but she may split
her budget across the datasets arbitrarily.
Otherwise, she spends as much money on her preferred dataset as possible,
and then as much money on the other dataset as possible.
Assume all poor buyers behave identically.

The sellers are, therefore, in a pricing game with each other.
Does an (approximately) stable solution exist? We show that the answer is `no'.
Formally, a pair $(p, q)$ of pricing strategies is a $c$-approximate Nash equilibrium ($c$-NE)
if no player can unilaterally deviate to a different pricing strategy
to increase her revenue by more than a factor of $c$.
We show that a $c$-NE doesn't exist for a large enough $c$.
